% AUTORIDAD DE FUENTE: Exposé IX §§13--13.2.2, líneas francesas
% 3952--4034; índices de escaneo 495--500, folios 484--489.

\medskip
\noindent
13. \underline{Relaciones con la teoría trascendente: el caso analítico
complejo}.

Para las nociones generales relativas a los espacios analíticos,
remitimos a [2].

\medskip
\noindent
13.1. \underline{Recordatorio sobre los grupos analíticos sobre $S$.}
Sea $S$ un espacio analítico y sea $A$ un espacio analítico sobre $S$
dotado de una estructura de grupo conmutativo sobre $S$ (es decir, $A$
es un grupo conmutativo en la categoría de los espacios analíticos sobre
$S$). Suponemos que $A$ es \underline{liso} sobre $S$ y denotamos por

%%%%%%%%%%  scan idx 496  |  folio 485  %%%%%%%%%%

\begin{equation}\tag{13.1.1}
 \underline E=\underline{\operatorname{Lie}}(A/S)
\end{equation}

\noindent
el módulo localmente libre sobre $S$ obtenido llevando el haz tangente
relativo de $A$ sobre $S$ a la sección cero. Admitiremos aquí la
definición del homomorfismo exponencial

\begin{equation}\tag{13.1.2}
 \exp:E\longrightarrow A\qquad ,
\end{equation}

\noindent
donde $E$ es el fibrado vectorial analítico sobre $S$ cuyo módulo de
secciones holomorfas es $\underline E$. En un entorno de un punto
$s\in S$, este homomorfismo se define primero cerca de la sección cero de
$E$ por los métodos familiares de la teoría de los grupos formales en
característica cero, y se prolonga después de manera evidente, por
multiplicatividad, a todo $E$. Es un homomorfismo de grupos analíticos
relativos sobre $S$; es \underline{étale}, y su imagen es el subespacio
abierto $A^o$ de $A$ que es la unión de las componentes neutras de las
fibras de $A$ sobre $S$. Si $R$ es el subespacio analítico núcleo de
$\exp$, tenemos, por tanto, una sucesión exacta de grupos analíticos
relativos

\begin{equation}\tag{13.1.3}
 0\longrightarrow R\longrightarrow E\longrightarrow A^o
 \longrightarrow0\qquad .
\end{equation}

\noindent
Aquí el morfismo $E\longrightarrow A^o$ es sobreyectivo y étale, por lo
que $R$ es un grupo analítico relativo \underline{étale} sobre $S$. Es
un subespacio analítico de $E$, y es cerrado en $E$ si y solo si $A^o$
es separado sobre $S$, es decir, si y solo si la sección cero de $A^o$
es cerrada en $A^o$.

\medskip
\noindent
13.1.4. Naturalmente, la sucesión exacta (13.1.3) es
\underline{funtorial en $A$} y conmuta con todo cambio de base
$S'\longrightarrow S$. Se obtiene así una equivalencia de categorías,
compatible con los cambios de base, entre la categoría de los grupos
analíticos relativos sobre $S$ que son lisos y tienen fibras conexas, y
la categoría de los pares $(E,R)$ formados por un fibrado vectorial
localmente libre $E$ sobre $S$ y un subgrupo analítico $R$ de $E$ que es
étale sobre $S$.

%%%%%%%%%%  scan idx 497  |  folio 486  %%%%%%%%%%

\medskip
\noindent
13.1.5. Cuando $A$ es un grupo «\underline{abeloide}» sobre $S$, con lo
que queremos decir que es liso, propio y separado sobre $S$, con fibras
conexas---de modo que, en particular, $A=A^o$---, entonces $R$ es un
grupo localmente constante sobre $S$, cuyo rango en cada punto $s$ (como
$\mathbb Z$-módulo) es igual a $2\operatorname{rank}_s(E)$, y
recíprocamente. Se obtiene así la interpretación familiar de un espacio
abeloide sobre $S$ en términos de un «\underline{sistema local}» $R$
sobre $S$ y de una estructura de fibrado vectorial analítico
\underline{complejo} sobre el fibrado analítico real
$R\otimes_{\mathbb Z}\mathbb R$ definido por $R$.

\medskip
\noindent
13.2. \underline{Completado formal a lo largo de una fibra.} Fijemos
ahora un punto $s\in S$ y trabajemos localmente en un entorno de $s$.
Para cada entero $n\geqslant0$, denotemos por $S_n$ el $n$-ésimo entorno
infinitesimal de $s$ en $S$ y, para todo espacio analítico $X$ sobre
$S$, pongamos asimismo $X_n=X\times_S S_n$. Al variar $n$, el sistema de
los $X_n$ forma entonces un sistema inductivo de espacios analíticos
sobre $S$, que denotamos por $\widehat X$ y llamamos el
\underline{completado formal de $X$ a lo largo de $X_s$}. Depende
claramente de manera funtorial de $X$, y su formación conmuta con los
límites proyectivos finitos; por tanto, $X\longmapsto\widehat X$ lleva
grupos a grupos. Obsérvese que $\widehat X$ depende, de hecho, únicamente
de la restricción de $X$ a un entorno abierto arbitrario $U$ de $s$ en
$S$.

%%%%%%%%%%  scan idx 498  |  folio 487  %%%%%%%%%%

\medskip
\noindent
\underline{Lema 13.2.1.} \underline{Sean los datos los de 13.1, y
supongamos que se cumplen las condiciones siguientes:} a) $A=A^o$;
b) $R$ \underline{es constante en un entorno de $s$}; y c) $R_s$
\underline{engendra el espacio vectorial complejo $E_s$}.
\underline{Sea $A'$ un segundo grupo analítico liso sobre $S$, y
consideremos el morfismo}

\begin{equation*}
 u\longmapsto\widehat u:
 \operatorname{Hom}_{S\text{-}\mathrm{gr}}(A,A')
 \longrightarrow
 \operatorname{Hom}_{\widehat S\text{-}\mathrm{gr}}
   (\widehat A,\widehat {A}')\quad ,
\end{equation*}

\noindent
\underline{así como el morfismo análogo obtenido restringiendo a un
entorno abierto variable $U$ de $s$ y pasando al límite inductivo:}

\begin{equation}\tag{13.2.1.1}
 u\longmapsto\widehat u:
 \operatorname{Hom}_{(S,s)\text{-}\mathrm{gr}}(A,A')
 \longrightarrow
 \operatorname{Hom}_{\widehat S\text{-}\mathrm{gr}}
   (\widehat A,\widehat {A}')\quad .
\end{equation}

\noindent
(\underline{Aquí $(S,s)$ denota el germen del espacio analítico $S$ en
$s$.}) \underline{Entonces este último morfismo es biyectivo.}

En efecto, consideremos la sucesión exacta

\begin{equation*}
 0\longrightarrow R'\longrightarrow E'\longrightarrow {A'}^o
 \longrightarrow\sisi{C}{0}
\end{equation*}

\noindent
análoga a (13.1.3) y definida por $A'$. Como los homomorfismos
$A\longrightarrow A'$ factorizan necesariamente por ${A'}^o$, y lo
mismo ocurre después de los cambios de base $U\longrightarrow S$ y
$S_n\longrightarrow S$, podemos suponer que $A'={A'}^o$. Por el
diccionario de 13.1.4, queda por demostrar que dar un germen de
homomorfismo de fibrados vectoriales $E\longrightarrow E'$ que lleve $R$
en $R'$ equivale a dar un sistema coherente de homomorfismos
$E_n\longrightarrow E'_n$ que lleven $R_n$ en $R'_n$.

%%%%%%%%%%  scan idx 499  |  folio 488  %%%%%%%%%%

Si $\underline{\mathcal O}_s$ denota el anillo local de $S$ en $s$, dar
un sistema coherente de homomorfismos $E_n\longrightarrow E'_n$ equivale
a dar un homomorfismo de $\widehat{\underline{\mathcal O}}_s$-módulos

\begin{equation*}
 \widehat F:\widehat{\underline E}_s
 \longrightarrow\widehat{\underline E}'_s,
\end{equation*}

\noindent
donde $\underline E_s$ y $\underline E'_s$ son los
$\underline{\mathcal O}_s$-módulos libres obtenidos en $s$ a partir de
los $\sisi{\underline{\mathcal O}_s}{\underline{\mathcal O}_S}$-módulos $\underline E$ y $\underline E'$,
correspondientes a $E$ y $E'$. Por otra parte, dar un germen de
homomorfismo de fibrados vectoriales $E\longrightarrow E'$ equivale a
dar un homomorfismo de $\underline{\mathcal O}_s$-módulos

\begin{equation*}
 F:\underline E_s\longrightarrow\underline E'_s\quad .
\end{equation*}

\noindent
Esto muestra ya que (13.2.1.1) es inyectivo. Para la sobreyectividad,
basta demostrar que, si $\widehat F$ es un homomorfismo como el anterior
tal que, para cada $n$, el homomorfismo inducido
$F_n:E_n\longrightarrow E'_n$ lleva $R_n$ en $R'_n$, entonces
$\widehat F$ procede de un $F$---es decir, lleva $\underline E_s$ en
$\underline E'_s$---y el germen correspondiente de homomorfismos
$E\longrightarrow E'$ lleva $R$ en $R'$.

%%%%%%%%%%  scan idx 500  |  folio 489  %%%%%%%%%%

Por la hipótesis c), $R_s=R_o$ engendra $E_s$; se sigue del lema de
Nakayama que el grupo $R_o$ (identificado bien con el grupo de secciones
de $\widehat R$ sobre $\widehat S$, bien con el grupo de gérmenes de
secciones de $R$ sobre $S$) genera $\sisi{\widehat E_s}{\widehat{\underline E}_s}$. Así, para demostrar
que $\widehat F(\underline E_s)\subset\underline E'_s$, basta demostrar
que $\widehat F(R_o)\subset\underline E'_s$. Pero
$\widehat F(R_o)\subset R'_o$, donde $R'_o$ se interpreta como el grupo
de gérmenes de secciones de $R'\subset E'$ en un entorno de $s$. Esto da
la inclusión afirmada y, por tanto, la existencia de $F$. Además, vemos
que $F(R_o)\subset R'_o$, donde $R_o$ y $R'_o$ se consideran como grupos
de gérmenes de secciones de $E$ y $E'$, respectivamente. Por b), $R_o$
genera $R$ en un entorno de $s$; se sigue que $F(R)\subset R'$ en un
entorno de $s$, lo que concluye la demostración.

\medskip
\noindent
\underline{Corolario 13.2.2.} \underline{Si $A$ satisface las
condiciones a), b) y c) de \sisi{13.2.2}{13.2.1}, queda determinado,
salvo isomorfismo único en un entorno de $s$, por el conocimiento del
$\widehat S$-grupo $\widehat A$.}
